\paragraph*{Objectif du chapitre :}
\begin{itemize}
\item  Vocabulaire et notation: puissance, exposant, ensemble des nombres décimaux, $\mathbb{D}$.fraction, nombre décimal, écriture décimale, nombre rationnel, ensemble des nombres rationnels, $\mathbb{Q}$.
\item  Multiplication, division et puissance de puissance.
\item  Additionner des fractions: mise au même dénominateur, multiplication et division de fractions.
\end{itemize}

\vspace{-3mm}
\section{Puissance d'un nombre}


%L'utilisation des puissances simplifie l'écriture des produits comportant le même facteur.
\noindent \begin{minipage}{12.5cm}
\begin{defi}{}{}[Puissance à exposant positif]
Si $x$ est un nombre et $n$ un entier naturel, alors $x^n$, qui se lit \og $x$ puissance $n$ \fg{} ou \og $x$ exposant $n$ \fg{}, est le nombre \[ x^n=\underset{n\ \text{fois} }{\underbrace{x\times \dots \times x}} \]
Par convention : pour tout $n \in \N$, $0^{n} = 0$; pour tout $n \in \Z$, $a^1 = a$ \quad et \quad $a^0 = 1$;
\end{defi}
\end{minipage}
\hspace{3mm}
\begin{minipage}{3.8cm}
\begin{ex}{}{}
\begin{align*}
2^{4} &= 2 \times 2 \times 2 \times 2 \\
&= 16\\
\\
(-3)^{3} &= (-3) \times (-3) \times (-3)\\
&=27
\end{align*}
\end{ex}
\end{minipage}

\noindent \begin{minipage}{11cm}
\begin{defi}{}{}[Puissance à exposant négatif]
Soit $n$ un nombre strictement positif et a un nombre relatif non nul. On appelle « $a$ puissance moins $n$ » le nombre noté $a^{-n}$ tel que :  $$a^{-n} = \dfrac{1}{a^n}$$
\end{defi}
\end{minipage}
\hspace{5mm}
\begin{minipage}{5cm}
\begin{ex}{}{}
$2^{-4} = \dfrac{1}{2^4} = \dfrac{1}{16}$ 

$(-3)^{-3} = \dfrac{1}{(-3)^{-3}} = \dfrac{1}{-27} = -\dfrac{1}{27}$  ;
\end{ex}  
\end{minipage}

\noindent \begin{minipage}{11cm}
\begin{prop}{}{}[Produit et quotient de puissances]
Soit $a$ un nombre relatif non nul; $m$ et $n$ deux entiers.
\begin{center}
$a^n \times a^m = a^{n+m}$ \hspace{0.5cm} et \hspace{0.5cm} $\dfrac{a^n}{a^m} = a^{n-m}$
\end{center}
\end{prop}
\end{minipage}
\hspace{5mm}
\begin{minipage}{5cm}
\begin{ex}{}{}
$3^5 \times 3^8 = 3^{5+8} = 3^{13}$

$\dfrac{5^2}{5^7} = 5^2 \times 5^{-7} = 5^{2-7} = 5^{-5}$
\end{ex} 
\end{minipage}

\noindent \begin{minipage}{11cm}
\begin{prop}{}{}[Puissance de produit et de quotient]
Soit $a$ et $b$ deux nombres relatifs non nuls, soit $n$ un entier.
\begin{center}
$a^n \times b^n = (a\times b)^{n}$ \hspace{0.5cm} et \hspace{0.5cm} $\dfrac{a^n}{b^n} = \left( \dfrac{a}{b} \right)^{n}$
\end{center}
\end{prop}
\end{minipage}
\hspace{5mm}
\begin{minipage}{5cm}
\begin{ex}{}{}
$3^5 \times 2^5 = (3\times 2)^{5} = 6^5$

$\dfrac{5^3}{10^3} = \left( \dfrac{5}{10} \right)^{3} = \left( \dfrac{1}{2} \right)^{3}$
\end{ex}
\end{minipage}

\noindent \begin{minipage}{11cm}
\begin{prop}{}{}[Puissance de puissance]
Soit $a$ un nombre relatif ; soit $m$ et $n$ des entiers
\begin{center}
$\left( a^n \right)^{m} = a^{n \times m}$
\end{center}
\end{prop}
\end{minipage}
\hspace{5mm}
\begin{minipage}{5cm}
\begin{ex}{}{}
$\left( 5^3 \right)^{4} = 5^{3 \times 4} = 5^{12}$

$\left( 2^4 \right)^{-5} = 2^{4 \times (-5)} = 2^{-20} = \dfrac{1}{2^{20}}$
\end{ex}
\end{minipage}

\begin{exec}
\hspace{0.5cm} \hypertarget{Ex_2_1_1}{} Exercice \ref{2_1_1} \hspace{0.5cm};
\hspace{0.5cm} \ding{79} \hypertarget{Ex_2_1_2}{} Exercice \ref{2_1_2}
\end{exec}

\section{Les nombres décimaux : $\mathbb{D}$}

\vspace{-2mm}
\begin{defi}{}{}
Un nombre est décimal si et seulement si il admet un développement décimal limité, c'est-à-dire si et seulement si il possède une écriture avec un nombre fini de chiffres après la virgule.

\noindent L'ensemble des nombres décimaux est noté $\mathbb{D}$. 
\end{defi}

\begin{rem}{}{}
\begin{itemize}
\item  Lorsqu'un nombre, dans son écriture décimale, n'a que des zéros à partir d'un certain nombre de décimales alors on dit que c'est un nombre décimal. 

Ainsi: $\np{2315,895643}= \np{2315,895643000}\dots$ est un nombre décimal.
	
Tous les nombres décimaux peuvent s'écrire sous forme de fraction: $\np{2315,895643}=\dfrac{\np{2315895643}}{\np{1000000}}$.
\item  Plus généralement, les nombres décimaux peuvent s'écrire comme quotient d'un entier relatif par une puissance d'exposant positif de 10 :

Un nombre est décimal s'il peut s'écrire $\dfrac{a}{10^n}$, $a$ appartenant à $\Z$ et $n$ à $\N$
\end{itemize}
\end{rem}

\bigbreak
Avec les définition ci-dessus, on observe plusieurs manières d'écrire un nombre :

\begin{itemize}
\item \textbf{l'écriture décimale} : La partie décimale compte un nombre fini de chiffres non nuls.

Par exemple, $237,698$ qui a pour partie entière $237$ et pour partie décimale $698$

\item \textbf{écriture scientifique} : C'est l'écriture sous forme du produit d'un nombre décimal ayant un seul chiffre non nul avant la virgule et d'une puissance de 10 (d'exposant entier relatif) 

Par exemple, $38 = 3,8 \times 10^1$ \hspace{0.8cm} ou \hspace{0.8cm} $0,0562 = 5,62 \times 10^{-2}$ \hspace{0.8cm}ou  encore \hspace{0.8cm} $-97631 = - 9,7631 \times 10^{4}$
\end{itemize}

\begin{exec}
\hspace{0.5cm} \hypertarget{Ex_2_1_3}{} Exercice \ref{2_1_3} \hspace{0.5cm}
\end{exec}

\section{Les fractions}

\begin{defi}{}{}
Une fraction est un nombre qui peut s'écrire comme le quotient d'un entier (relatif) par un autre entier (relatif non nul).
\end{defi}
\medskip

\noindent Au lycée et ultérieurement on utilise que la notation fractionnaire $\left(\dfrac{1}{3}\right)$ pour la division et plus le $\div$ $(1 \div 3)$.
\medskip

\begin{rem}{}{}
\begin{itemize}
\item Il existe des nombres qui ne peuvent pas s'écrire sous forme de fractions: $\sqrt{2}$, $\pi$ (admis).
\item Rappelons une règle de priorité opératoire: dans une fraction la division correspondant à la fraction est la dernière opération effectuée.
\item La manipulation des fractions nécessite de connaître les règles :

\begin{minipage}{7cm}
\begin{center}
 d'addition (réduire au même dénominateur),
 \end{center} \[  \frac{a}{b}+\frac{x}{y}=\frac{ay+xb}{by} \] \end{minipage}
\vline
\begin{minipage}{5cm}
\begin{center}
 de multiplication,
 \end{center} \[  \frac{a}{b}\times \frac{x}{y}=\frac{a \times x}{b \times y} \] 
\end{minipage}
\vline
\begin{minipage}{4cm}
\begin{center}
de division 
\end{center}\[  \frac{ \frac{a}{b} }{ \frac{x}{y}}=\frac{a}{b}\times \frac{y}{x}. \] 
\end{minipage}

\medskip
Rappelons encore que pour simplifier une fraction il faut qu'il y ait un facteur \\commun au numérateur et au dénominateur: \[  \frac{a\times {\color{orange}x}}{b\times {\color{orange}x}}=\frac{a}{b} \]
\item Une astuce d'écriture qui nous sera souvent utile \[ \frac{a}{b}=\frac{1}{b}\times a \]
\item Rappelons enfin qu'il est impossible de diviser par $0$. 

\textbf{Raisonnons par l'absurde: }

Supposons qu'il existe un nombre $q$ tel que $\frac{3}{0}=q$ (nous supposons que l'on peut diviser $3$ par $0$). 

Donc $3=0\times q$. 

Ce qui est impossible car $0\times q=0$ ($0$ est une valeur absorbante). 

Nous avons donc démontré par l'absurde qu'il est impossible de diviser $3$ par $0$.
\end{itemize}
\end{rem}

\begin{prop}{}{}
\noindent Les règles de priorités pour l'instant se limitent à:
\begin{itemize}
\item Les calculs entre parenthèses ou crochets. S'il y a des parenthèses emboîtées les plus emboîtées sont prioritaires. La barre de fraction joue le même rôle que des parenthèses autour des numérateurs et dénominateurs.
\item Les exposants (puissances)
\item Les multiplications et division ($\div$) en allant de gauche à droite.
\item Les additions et soustractions en allant de gauche à droite.
\end{itemize}
\end{prop}	

\begin{exec}
\hspace{0.5cm} \hypertarget{Ex_2_2_1}{} Exercice \ref{2_2_1} \hspace{0.5cm};
\hspace{0.5cm} \hypertarget{Ex_2_2_2}{} Exercice \ref{2_2_2} \hspace{0.5cm};
\hspace{0.5cm} \hypertarget{Ex_2_2_3}{} \ding{79} Exercice \ref{2_2_3} \hspace{0.5cm};
\hspace{0.5cm} \hypertarget{Ex_2_2_4}{} Exercice \ref{2_2_4} \hspace{0.5cm};
\hspace{0.5cm} \hypertarget{Ex_2_2_5}{} \ding{79} Exercice \ref{2_2_5} \hspace{0.5cm}
\end{exec}

\vspace{-2mm}
\section{Les nombres rationnels : $\mathbb{Q}$}

\vspace{-2mm}
\begin{defi}{}{}
Les nombres rationnels sont les nombres qui peuvent s'écrire sous forme de fraction. 

\noindent L'ensemble de tous les nombres rationnels est noté: $\mathbb{Q}$.
\end{defi}

\begin{rem}{}{}
\begin{itemize}
%\item  Un entier est un nombre rationnel. Par exemple: $-13=\dfrac{-13}{1}$.
%\item  Les nombres décimaux peuvent s'écrire comme des fractions donc ce sont des nombres rationnels. \\ Ainsi: $0,1=\dfrac{1}{10}$.
\item  Il existe des nombres qui ne sont pas décimaux et qui sont rationnels. \\ Ainsi $\dfrac{1}{3}$ est un nombre rationnel mais son écriture décimale est infinie donc ce n'est un nombre décimal.
\begin{deme}{}{}
Proposition à démontrer : \og $\dfrac{1}{3}$ n'est pas décimal \fg.

\medskip
On va démontrer cette proposition par l'absurde :

\noindent Supposons que $\dfrac{1}{3}$ soit décimal. Alors il est de la forme $\dfrac{a}{10^n}$, $a$ appartenant à $\Z$ et $n$ à $\N$.

\noindent On en déduit que : 
\begin{align*}
\dfrac{1}{3} = \dfrac{a}{10^n} \quad \Longleftrightarrow \quad
\dfrac{1}{3} {\color{cocours} \times 3}= \dfrac{a}{10^n}{\color{cocours} \times 3}\quad \Longleftrightarrow \quad
1 = \dfrac{3a}{10^n}\quad \Longleftrightarrow \quad
1 {\color{cocours} \times 10^n}= \dfrac{3a}{10^n}{\color{cocours} \times 10^n}\quad \Longleftrightarrow \quad
10^n = 3 \times a
\end{align*}

\noindent Alors l'une des puissances de 10 serait un multiple de 3. 

\noindent C'est faux d'après les critères de divisibilité (et sans même avoir à poser la division décimale).

\noindent Donc $\dfrac{1}{3}$ n'est pas décimal.\hfill $\blacksquare$
\end{deme}
\item  Tous les nombres ne sont pas rationnels.
\item  Lorsqu'un nombre rationnel a une écriture décimale infinie on préfère son écriture fractionnaire.
\item  Un nombre rationnel admet une écriture fractionnaire canonique appelée la forme irréductible qui est obtenue lorsque le PGCD des numérateur et dénominateur est 1 (ils n'ont pas de facteur commun).
	
Certaines calculatrices donnent la forme irréductible.
\end{itemize}
\end{rem}

\begin{prop}{}{}
Lorsqu'un nombre a une écriture décimale infinie périodique c'est un nombre rationnel.
\end{prop}

\begin{rem}{}{}
On dit que l'écriture est périodique lorsqu'une série de chiffre se reproduit indéfiniment.
\end{rem}

\begin{ex}{}{}
\begin{itemize}
	\item  $0,333\dots$ est rationnel, en effet c'est l'écriture décimal de la fraction $\dfrac{1}{3}$
	\item  $45,78242424\dots$ est aussi rationnel, bien que sa forme fractionnaire sois plus dur à déterminer
\end{itemize}
\end{ex}